\documentclass[11pt]{amsart}
\usepackage{geometry}                % See geometry.pdf to learn the layout options. There are lots.
\geometry{letterpaper}                   % ... or a4paper or a5paper or ... 
%\geometry{landscape}                % Activate for for rotated page geometry
%\usepackage[parfill]{parskip}    % Activate to begin paragraphs with an empty line rather than an indent
\usepackage{graphicx}
\usepackage{amssymb}
\usepackage{epstopdf}
\DeclareGraphicsRule{.tif}{png}{.png}{`convert #1 `dirname #1`/`basename #1 .tif`.png}

\title{PS 3.30}
\author{Enrique Zuñiga Zuñiga}
%\date{}                                           % Activate to display a given date or no date

\begin{document}
\maketitle
%\section{}
\noindent Es posible demostrar que la suma de un número “suficiente” de términos de la serie:\newline\newline
$$Sen(x) = x-\tfrac{x^3}{3!} +\tfrac{x^5}{5!}-\tfrac{x^7}{7!}+ . . . $$\newline\newline
Es un número considerablemente cercano a SEN(X), y que la diferencia entre SEN(X) y la suma antes mencionada se vuelve menor conforme se toman más y más términos. Escriba un diagrama de flujo tal que dada una X cualquiera calcule el SEN(X) utilizando la serie anterior, de tal modo que la diferencia entre la serie y un nuevo término sea menor o igual a 0.01. Imprima el número de términos requerido para obtener esta precisión.\newline\newline
Dato: X (variable de tipo entero que representa el número que se ingresa).\newline
%\subsection{}



\end{document}  